% Appendix, uploaded SEPARATELY from the main paper (MLSys CFP). Build: make appendix
% Reviewers are not required to read it: nothing the main paper depends on goes here.
% Main-paper references (Draft 3): Appendix A, B, C. Plan from GDoc Draft 3:
%   A: lr-1e-3 decomposition (hierarchy ~70% + edge increment, BH-adjusted), Top-K tables at lr 1e-3
%      and tuned, inception_v3 per-graph check, FIG-R2 (xla:random paired seeds), FIG-C2 (roofline).
%   B: nlp:default amendment table + pre-registration breach disclosure (what was read, and when).
%   C: protocol notes and robustness: every headline contrast at eval batch 8 and 32
%      (appendix_eval_batch.tex), capstone_robustness.tex, checkpoint selection, CPU-trained arms (O17),
%      run-to-run variation (O22), oracle epochs (O23), LR scope, CONV subset size (O20), FIG-M2
%      stability curve (batch 8: median SD 0.0145 at 7 graphs, 0.0131 at 20; batch 32: 0.0152 / 0.0135),
%      hyperparameters.tex.
%   D: audit rubric, per-paper scores of the 20, and Frame B Axis 3 (O15; experiments/axis3_recent_coverage_scores.md): B8-B10 author clustering, N/A count, "most-cited in this frame",
%      the declaration's blindness statement, Amendment 1 (4c400ff; cite the anonymized hash).
\documentclass{article}
\usepackage{microtype}
\usepackage{graphicx}
\usepackage{booktabs}
\usepackage{amsmath}
\usepackage{xcolor}
\usepackage{hyperref}
\newcommand{\theHalgorithm}{\arabic{algorithm}}
\usepackage{mlsys2025}
\newcommand{\TODO}[1]{\textcolor{red}{[TODO: #1]}}
\newcommand{\SLOT}[1]{\textcolor{blue}{[SLOT: #1]}}

\mlsystitlerunning{Appendix}

\begin{document}
\twocolumn[
\mlsystitle{Appendix: What Do Static Computation Graphs Know About Hardware Performance?}
\begin{mlsysauthorlist}
\mlsysauthor{Anonymous Author}{anon}
\end{mlsysauthorlist}
\mlsysaffiliation{anon}{Anonymous Institution}
\mlsyscorrespondingauthor{Anonymous Author}{anon@example.com}
\vskip 0.3in
]
\printAffiliationsAndNotice{}

\appendix
% Section order follows the main text's references: A (Result I), B (nlp:default + pre-registrations),
% C (protocol and robustness), D (audit).
\section{Result I: decomposition and per-graph checks}
\label{app:result1}
% lr-1e-3 decomposition (hierarchy ~70% + edge increment, BH), Top-K tables, inception_v3, FIG-R2, FIG-C2.
\TODO{Result I extras}

\section{Pre-registrations, amendments and the nlp:default failure}
\label{app:prereg}
% 6d (+ Amendments 1 and 2), 6e, 8c, 8d: what was registered, when (commit), and outcome.
% 6d breach disclosure: what nlp:default tau was read, when, before Amendment 1.
% Constant-output collapse, repro, lr diagnosis.
\TODO{preregistrations and nlp:default}

\section{Protocol details and robustness}
\label{app:robustness}
% Hyperparameter table (tables/hyperparameters.tex), the five pipeline bug fixes, compute harness.
% Every headline contrast at eval batch 8 and 32 (tables/appendix_eval_batch.tex),
% tables/capstone_robustness.tex, O9 Top-K, 6e flat LR curve, FIG-M2, O17/O20/O22/O23.
\TODO{protocol and robustness}

\section{Audit: rubric, scores and coverage}
\label{app:audit}
% Facts from experiments/meta_analysis_rubric.md, phase7_meta_analysis.md, phase7_coverage/screening.md,
% axis3_recent_coverage_declaration.md and axis3_recent_coverage_scores.md. Tables are generated by
% scripts/paper_figures.py. Prose drafted by Claude at the first author's request: revise as needed.

\subsection{Rubric}
Each paper's headline comparative architectural claim is scored on five axes (Table~\ref{tab:rubric}).
The rubric was fixed before any paper was scored and applied identically to every paper. When the
evidence is ambiguous, the weaker category is scored. Every score is backed by a quoted passage or
an explicit ``not stated''.

\begin{table*}[t]
\caption{Audit rubric. Each axis is scored pass (P), partial ($\partial$) or fail (F) on the paper's
headline comparative claim.}
\label{tab:rubric}
\centering
\small
\begin{tabular}{p{2.6cm}p{4.6cm}p{4.6cm}p{4.2cm}}
\toprule
Axis & P & $\partial$ & F \\
\midrule
1. Seeds / variance & $\geq$3 seeds, mean and dispersion reported & several runs but no dispersion, best-of-$N$
  only, or seeds for results other than the headline & single run or unstated seeds \\
2. Parameter matching & explicit parameter or FLOP match within a stated tolerance, or a capacity sweep
  & similar size but no explicit match, or parameters matched but not training budget & proposed model
  larger or longer-trained with no control \\
3. Significance testing & a named test or statistical CI on the comparison, with statistic or $p$-value
  & variance reported but no test & no variance and no test \\
4. Evaluation-set integrity & held-out test set of adequate size, selection on a separate split & honestly
  labelled validation-only, or a very small test set without a power caveat & reported on the selection
  set, or split never stated \\
5. Effect vs noise and baseline & gain clearly exceeds the variance and the baseline is tuned and
  competitive & gain comparable to its error bars, or a lightly tuned baseline & sub-noise gain, or an
  under-tuned baseline \\
\bottomrule
\end{tabular}
\end{table*}

\subsection{Selection and verification}
The 20 papers (2018--2026) are a purposive sample: two pilot papers (the benchmark we use and one
architecture paper on it), eight papers from the established lineage of learned cost and performance
models, and ten found by web searches whose terms, sources and dates were not logged. The selection
protocol was written after scoring. The second round was steered towards finding counter-examples,
which biases the sample towards papers that pass axes. Candidates that were excluded or never scored
are listed in the released repository.

All 20 papers were first scored by an AI model from the papers' text. Five papers (A5, A11, A14, A16
and A18; 25\% of the sample, chosen at random) were then verified by a human on all five axes, which
corrected 3 of 25 scores, all towards stricter. Axis~3 was then human-verified for all 20 against each
paper's PDF; one score became looser (A2, F to $\partial$: it reports $\pm1\sigma$ but runs no test).
Axes 1, 2, 4 and 5 of the other 15 papers are AI-prescreened only. A separate AI rater scored the same
five papers blind as a pilot (21 of 25 scores agreed; Cohen's $\kappa \approx 0.67$); this is not an
independent human second rater.

\paragraph{Coverage.} Before running it, we declared a citation-ranked coverage check: six fixed Semantic
Scholar queries (2018--2026, top 100 results each), retrieved on 2026-09-27, returned 571 candidates.
Titles and abstracts were screened against the scope rule by an AI model, with ambiguous abstracts
screened in and all 31 boundary decisions checked by one of the authors. Our sample contains 8 of the
30 most-cited in-scope papers (Frame A) and 0 of the 10 most-cited in-scope papers from 2024--2026
(Frame B). The queries retrieve only 10 of our 20 papers. In-scope papers outside the sample were not
audited, except for the Axis~3 check below.

\begin{table*}[t]
\caption{Per-paper audit scores for the 20 papers. ``Verified'' gives the axes checked by a human
against the paper; the other axes are AI-prescreened.}
\label{tab:audit-per-paper}
\centering
\small
% Generated by scripts/paper_figures.py from the per-paper tables in experiments/phase7_meta_analysis.md
\begin{tabular}{rlcccccc}
\toprule
\# & Paper & Ax.\,1 & Ax.\,2 & Ax.\,3 & Ax.\,4 & Ax.\,5 & Verified \\
\midrule
A1 & TPUGraphs~\cite{phothilimthana2023tpugraphs} & F & F & F & $\partial$ & F & Axis 3 \\
A2 & TGraph~\cite{khizbullin2024graph} & $\partial$ & F & $\partial$ & F & $\partial$ & Axis 3 \\
A3 & GRANITE~\cite{sykora2022granite} & F & F & F & P & $\partial$ & Axis 3 \\
A4 & Ithemal~\cite{mendis2019ithemal} & F & F & F & P & $\partial$ & Axis 3 \\
A5 & ProGraML~\cite{cummins2021programl} & F & F & F & $\partial$ & $\partial$ & all axes \\
A6 & TLP~\cite{zhai2022deep} & F & F & F & P & $\partial$ & Axis 3 \\
A7 & PerfSAGE~\cite{chai2023perfsage} & F & F & F & P & F & Axis 3 \\
A8 & NeuSight~\cite{lee2024forecasting} & F & F & F & P & $\partial$ & Axis 3 \\
A9 & Kaufman et al.~\cite{kaufman2021learned} & $\partial$ & $\partial$ & F & $\partial$ & $\partial$ & Axis 3 \\
A10 & AutoTVM~\cite{chen2018learning} & F & F & F & $\partial$ & $\partial$ & Axis 3 \\
A11 & TenSet~\cite{zheng2021tenset} & F & F & F & $\partial$ & F & all axes \\
A12 & Polynormer~\cite{deng2024polynormer} & F & F & F & F & $\partial$ & Axis 3 \\
A13 & Guha et al.~\cite{guha2025learned} & $\partial$ & F & F & F & $\partial$ & Axis 3 \\
A14 & Facile~\cite{abel2023facile} & F & F & F & $\partial$ & $\partial$ & all axes \\
A15 & COGNATE~\cite{sudusinghe2025cognate} & F & F & F & P & $\partial$ & Axis 3 \\
A16 & RLM for code~\cite{akhauri2025regression} & F & F & F & F & F & all axes \\
A17 & Omniwise~\cite{wang2025omniwise} & F & $\partial$ & F & P & F & Axis 3 \\
A18 & HLS-GNN~\cite{wu2022highlevel} & $\partial$ & F & F & P & F & all axes \\
A19 & Compiler world models~\cite{pan2026toward} & F & F & F & $\partial$ & $\partial$ & Axis 3 \\
A20 & PerfVec~\cite{li2024learning} & F & $\partial$ & F & P & $\partial$ & Axis 3 \\
\bottomrule
\end{tabular}

\end{table*}

\subsection{Axis 3 check on recent work (Frame B)}
\label{app:frameb}
Because the sample covers none of Frame B, we checked those ten papers on Axis~3 alone. The check
was declared and frozen before any of the ten papers was opened in full (commit \texttt{e12c51a},
2026-10-12, in the released repository). The declaration fixes the paper list (no substitutions),
the rule for choosing each paper's headline comparison, the scoring anchors and the sentence
reported in Section~8 for each possible outcome.

\paragraph{Headline comparison.} The scored comparison is between the paper's cost or performance
model and a baseline model where only the model differs; comparisons against frameworks, vendor
libraries or compilers, or comparisons that also change the search, search space or budget, do not
qualify, and neither does a comparison where the paper is silent on whether anything else differs.
The headline is the first qualifying comparison in the abstract, else in the Introduction's
contribution list, else in the evaluation section. It was recorded verbatim before the paper's
results were examined. A paper with no qualifying comparison would be scored N/A.

\paragraph{Scoring.} One author scored all ten papers from the version of record where obtainable,
else the latest arXiv version, using the main PDF and any appendix inside it. No AI pre-screened
the papers, extracted quotes or suggested scores. The scoring window was originally 2026-10-19 to
2026-10-22; Amendment~1 (\SLOT{anonymized hash of Amendment 1}, 2026-10-28, made before any listed
paper was opened) moved it to 2026-10-28 to 2026-10-29 because of the writing schedule, and
changed nothing else. All ten were scored on 2026-10-28. Time spent per paper was not recorded.

\paragraph{Result.} All ten qualify and all ten score F: none runs a significance test, and none
reports variance, for its headline comparison (Table~\ref{tab:frameb}). No paper was N/A, not
obtained or not scored, so $n = 10$ and $k = 0$. These results are reported separately and are never
combined with the 20-paper sample.

\paragraph{Caveats.} ``Most-cited'' means most-cited \emph{in this frame}: the ten papers had 2 to 16
citations at retrieval, so they are recent papers ranked by a noisy signal, not the field's landmark
papers. B8, B9 and B10 come from one research group and share 6--8 authors; B9 may be the journal
extension of B8. They were scored separately, as declared. Years are as listed by Semantic Scholar
at retrieval; for B1 the version of record is the PACT 2025 proceedings, and B9 and B10 appeared
online in 2025.

\paragraph{Blindness statement (from the frozen declaration, verbatim).}
\begin{quote}
\small
The scorer had not read the full text of any of the 10 papers before this declaration was frozen.
During coverage screening (2026-10-05) the scorer saw each paper's title and, where Semantic Scholar
had one, its abstract. None of the abstracts reports seeds, error bars, dispersion, confidence
intervals or a hypothesis test. Several use ``significantly'' or ``significant improvements'' in a
non-statistical sense, and these words were seen. One paper (NLTSP) had no abstract in Semantic
Scholar. The scorer may know some of these papers from general reading; that familiarity cannot be
ruled out, only the absence of a full-text read before freezing. Axis 3 was scored under anchors and
tie-break rules fixed before the 20-paper audit and unchanged since.
\end{quote}

\begin{table*}[t]
\caption{Axis~3 check of the ten most-cited in-scope 2024--2026 papers in the coverage frame (Frame B).
``Cit.'' is the Semantic Scholar citation count at retrieval (2026-09-27). Evidence gives where the
headline comparison's numbers are reported; none states variance or a test.}
\label{tab:frameb}
\centering
\scriptsize
% Generated by scripts/paper_figures.py from experiments/axis3_recent_coverage_scores.md (scores),
% the frozen declaration (list) and experiments/phase7_coverage/screening.md (citations at retrieval)
\begin{tabular}{rp{6.2cm}lrrcl}
\toprule
\# & Title & Venue & Year & Cit. & Axis 3 & Evidence \\
\midrule
B1 & LOOPer: A Learned Automatic Code Optimizer For Polyhedral Compilers & PACT & 2024 & 16 & F & not stated; \S{}VI-B4, p. 10 \\
B2 & Learning to Compare Hardware Designs for High-Level Synthesis & MLCAD & 2024 & 10 & F & not stated; Table 1, p. 5 \\
B3 & Pruner: A Draft-then-Verify Exploration Mechanism to Accelerate Tensor Program Tuning & ASPLOS & 2024 & 9 & F & not stated; Table 11, p. 14 \\
B4 & Fasor: A Fast Tensor Program Optimization Framework for Efficient DNN Deployment & ICS & 2024 & 5 & F & not stated; Tab. 8, p. 507 \\
B5 & Crop: An Analytical Cost Model for Cross-Platform Performance Prediction of Tensor Programs & DAC & 2024 & 4 & F & not stated; Table 2, p. 6 \\
B6 & Fast and Accurate Context-Aware Basic Block Timing Prediction using Transformers & CC & 2024 & 4 & F & not stated; Table 5, p. 234 \\
B7 & NLTSP: A cost model for tensor program tuning using nested loop trees & J. Syst. Archit. & 2025 & 3 & F & not stated; \S{}4.2.1 / Fig. 7, p. 8 \\
B8 & UniCoMo: A Unified Learning-Based Cost Model for Tensorized Program Tuning & ICCD & 2024 & 3 & F & not stated; Table IV, p. 493 \\
B9 & UniCoX: A Unified Cost Model for Tensorized Program Tuning Across Ubiquitous Accelerators & IEEE TC & 2026 & 2 & F & not stated; Fig. 17, p. 317 \\
B10 & UniSparTa: A Unified Sparse Tensor Program Tuning Framework & IEEE TCAD & 2026 & 2 & F & not stated; Fig. 13, p. 3370 \\
\bottomrule
\end{tabular}

\end{table*}

\section{Additional figures}
\label{app:figures}
% [[FIG-A1 | HAVE: phase4_topk.png]]  [[FIG-A2 | HAVE: diffpool_S_comparison_multi.png]]
% [[FIG-A3 | HAVE: training_curves.png]] (caption: cosine-schedule artifact)
\TODO{figures}

\bibliography{refs}
\bibliographystyle{mlsys2025}
\end{document}
